% from ega4-15-fr.tex
\oldpage[IV-3]{231}
équidimensionnel au point $x$ (15.4.1), d'où l'existence d'un morphisme
quasi-fini $g:U\to Y'$ tel que toute composante irréductible de $U$ domine $Y'$
(13.3.1, $b$)) et que $g_y:U_y\to Y_y'$ soit fini et surjectif (13.3.1.1).
D'autre part, pour prouver que $\mathcal F|U$ est $g$-plat au point $x$, il
suffit, puisque $h:Y'\to Y$ est plat, de montrer (en vertu du critère de
platitude par fibres (11.3.10)) que $\mathcal F_y$ est $g_y$-plat au point $x$.
Mais par hypothèse $\mathcal F_y$ est un $\mathcal O_{X_y}$-Module de
Cohen-Macaulay au point $x$ et $Y_y'$ un préschéma \emph{régulier}~; d'autre
part, comme $g_y$ est fini et surjectif, il vérifie la condition $e')$ de
(15.4.2) en vertu de (5.6.10)~; il suffit donc pour conclure d'appliquer (15.4.2)
à $g_y$ et à $\mathcal F_y$.
\end{proof}

% IV3_B32_P231_END_BEFORE_SUBSECTION_15.5

\subsection{Rang séparable des fibres d'un morphisme quasi-fini et universellement ouvert. Application aux composantes connexes géométriques des fibres d'un morphisme propre}
\label{subsection:IV.15.5-fr}

\begin{proposition}[15.5.1]
\label{IV.15.5.1-fr}
Soient $Y$ un préschéma localement noethérien, $f:X\to Y$ un morphisme séparé et quasi-fini. Pour tout $z\in Y$, soit $n(z)$ le nombre géométrique des points de $f^{-1}(z)$ (ou rang séparable de $f^{-1}(z)$ sur $\mathbf{k}(z)$, cf. \textbf{I}, 6.4.8). On a les propriétés suivantes~:
\begin{enumerate}[label=\emph{(\roman*)}]
  \item La fonction $z\mapsto n(z)$ est semi-continue inférieurement en tout point $y\in Y$ tel que $f$ soit universellement ouvert aux points de $f^{-1}(y)$.
  \item Soit $y$ un point de $Y$ tel que $f$ soit universellement ouvert aux points de $f^{-1}(y)$~; si la fonction $z\mapsto n(z)$ est constante dans un voisinage de $y$, alors $f$ est propre au point $y$ (15.7.1).
  \item On suppose que $f$ soit universellement ouvert et que tout point de $Y$ soit géométriquement unibranche~; alors, si la fonction $z\mapsto n(z)$ est localement constante dans $Y$, les composantes irréductibles de $X$ sont deux à deux disjointes.
\end{enumerate}
\end{proposition}

\begin{proof}
(i) Notons que puisque $f$ est quasi-fini, le nombre $n(z)$ est le nombre géométrique des composantes connexes de $f^{-1}(z)$~; on sait que la fonction $z\mapsto n(z)$ est localement constructible (9.7.9)~; compte tenu de ($\mathbf{0}_{\mathrm{III}}$, 9.3.4), on est ramené à prouver le

\begin{corollary}[15.5.2]
\EGAReaderTarget{ega4.c04.cor.15-5-2}{IV.15.5.2-fr}
Sous les hypothèses générales de (15.5.1), soit $y$ un point de $Y$ tel que $f$ soit universellement ouvert aux points de $f^{-1}(y)$. Alors, si $y'$ est une générisation de $y$, on a
\begin{equation}
\tag{15.5.2.1}
\label{IV.15.5.2.1-fr}
  n(y)\leqslant n(y').
\end{equation}
\end{corollary}

Pour tout changement de base $g:Z\to Y$, $f_{(Z)}:X_{(Z)}\to Z$ est séparé, quasi-fini et universellement ouvert aux points de $X_{(Z)}$ se projetant dans $f^{-1}(y)$~; en outre, si $z,z'$ sont deux points de $Z$ tels que $g(z)=y$, $g(z')=y'$ et que $z'$ soit une générisation de $z$, on a $n(z)=n(y)$ et $n(z')=n(y')$ (I, 6.4.12)~; il suffit donc de démontrer (15.5.2.1) pour un $g$ convenable. On peut évidemment supposer $y\ne y'$. Nous utiliserons le lemme suivant~:

\begin{lemma}[15.5.2.2]
\label{IV.15.5.2.2-fr}
Sous les hypothèses générales de (15.5.1), si $y$ est un point de $Y$, $y'$ une générisation de $y$ distincte de $y$, il existe un spectre d'anneau de valuation discrète $Z$, de point fermé $z$ et de point générique $z'$, et un morphisme $g:Z\to Y$ tels que~:
\begin{enumerate}[label=\emph{\arabic*\textsuperscript{\textdegree}}]
  \item $g(z)=y$, $g(z')=y'$~;
  \item si l'on pose $f'=f_{(Z)}$, $n(y)$ est le nombre de points de $f'^{-1}(z)$ et $n(y')$ est le nombre de points de $f'^{-1}(z')$.
\end{enumerate}
\end{lemma}

En effet, il existe un anneau de valuation discrète $A_1$ et un morphisme $g_1$ de $Z_1=\operatorname{Spec}(A_1)$ dans $Y$ tels que, si $z_1$ et $z'_1$ sont le point fermé et le point générique de $Z_1$,
